% AY2018 材料力学 〔II〕（転記）
% 出典: 04_材料力学/original/04_材料力学/2018/2018za.pdf
% 要約: はり\textcircled{1}, はり\textcircled{2}からなる片持ちはりに単位長さ当たり f0 の等分布荷重。
%       断面は長方形, 高さはそれぞれ 2h（はり\textcircled{1}）, h（はり\textcircled{2}）, 幅b。長さは共にL。
%       固定端A → はり\textcircled{1}(A-B, 高さ2h) → はり\textcircled{2}(B-C, 高さh) → 自由端C。縦弾性係数E。
%       (1) SFD/BMD, (2) I1,I2, (3) 最大せん断応力τ1,τ2, (4) 先端Cのたわみδ。

\section*{〔II〕}

図2に示すように, はり\textcircled{1}, はり\textcircled{2}からなる片持ちはりに単位長さ当たり $f_0$ の等分布荷重が
作用している. はり\textcircled{1}, はり\textcircled{2}の断面は長方形であり, 高さはそれぞれ $2h$ および $h$,
幅（図の奥行き方向）は $b$, はり\textcircled{1}, \textcircled{2}の長さは共に $L$ とする.
縦弾性係数（ヤング率）を $E$ とする. 次の問い (1)\,$\sim$\,(4) に答えよ.

\begin{center}
\begin{tikzpicture}[>=Latex,scale=0.95]
  % 固定壁（左, A）
  \fill[pattern=north east lines] (-0.4,-1.1) rectangle (0,1.6);
  \draw (0,-1.1) -- (0,1.6);
  % はり\textcircled{1}（A-B, 高さ2h: y -1.0..1.0）
  \draw (0,1.0) rectangle (4,-1.0);
  % はり\textcircled{2}（B-C, 高さh: y -0.5..0.5）
  \draw (4,0.5) rectangle (8,-0.5);
  % 中立軸
  \draw[dash dot] (-0.3,0) -- (8.7,0) node[right]{$x$};
  % y軸
  \draw[->] (0,0) -- (0,2.0) node[left]{$y$};
  % 等分布荷重 f0（全長, 上向き面に下向き矢印）
  \foreach \x in {0.4,1.0,...,3.8} {\draw[->] (\x,1.7) -- (\x,1.05);}
  \foreach \x in {4.2,4.8,...,7.8} {\draw[->] (\x,1.7) -- (\x,0.55);}
  \node at (5.0,2.0) {$f_0$};
  % 節点
  \node[below] at (0,-1.05) {A}; \node[below] at (4,-1.05) {B}; \node[below] at (8,-0.55) {C};
  % はりラベル
  \node at (1.3,0.55) {はり\textcircled{1}}; \node at (5.3,0.25) {はり\textcircled{2}};
  % 高さ 2h, h
  \draw[<->] (3.6,-1.0) -- (3.6,1.0) node[midway,fill=white]{$2h$};
  \draw[<->] (7.0,-0.5) -- (7.0,0.5) node[midway,fill=white]{$h$};
  % 長さ L, L
  \draw[<->] (0,-1.35) -- (4,-1.35) node[midway,fill=white]{$L$};
  \draw[<->] (4,-1.35) -- (8,-1.35) node[midway,fill=white]{$L$};
\end{tikzpicture}

\vspace{1mm}
図2
\end{center}

\begin{enumerate}[label=(\arabic*)]
  \item せん断力図（SFD）および曲げモーメント図（BMD）を求めよ.

  \item はり\textcircled{1}およびはり\textcircled{2}の断面二次モーメント $I_1$, $I_2$ を求めよ.

  \item はり\textcircled{1}およびはり\textcircled{2}における最大せん断応力 $\tau_1$, $\tau_2$ を求めよ.

  \item はり先端 C のたわみ $\delta$ を求めよ.
\end{enumerate}
